\chapter{DAGonFS: An In-Memory Scratch File System}
\label{app:dagonfs}

Every DAGonStar task writes into a folder under the workflow's scratch
directory, and every \texttt{workflow://} reference reads from one. These
files are short-lived. They are produced, read by a few consumers and
discarded. On a cluster the scratch directory usually lies on a shared
persistent file system such as NFS or Lustre. Every small file then pays the
metadata and durability costs of storage meant to last. DAGonFS (DAGon File
System) keeps the scratch space in the aggregate memory of the job's MPI
processes instead. It is the companion project of DAGonStar, and this
appendix accompanies course deck~A1.

\section{What DAGonFS Is}

DAGonFS is an \emph{ad-hoc, in-memory, MPI-distributed} file system exposed
through FUSE. One executable, \texttt{dagonfs}, mounts a directory whose file
contents are spread across the RAM of a set of MPI ranks. It provides no
durable storage: when the volume is unmounted, its contents are gone.

\begin{table}[h]
\centering
\begin{tabular}{@{}lll@{}}
\toprule
Model & Mount points & Who serves requests \\
\midrule
client-server (default) & one, on rank 0 & rank 0 holds the namespace;
  ranks $1\ldots P-1$ store blocks \\
peer-to-peer & \texttt{<root>/<rank>} & every peer serves its mount and
  replicates namespace changes \\
\bottomrule
\end{tabular}
\caption{The two deployment models of DAGonFS, chosen per volume at run time.}
\label{tab:dagonfs-models}
\end{table}

The model is a run-time property, selected with \texttt{--model}, the
\texttt{dagonfs\_model} key of \texttt{DAGonFS.ini}, or the C++ and Python
volume APIs. Use client-server when every task runs on the node of rank~0,
and peer-to-peer when tasks run on several nodes, each using the mount point
of a rank on its own node.

\section{Architecture}

A task's POSIX call enters the Linux VFS and reaches DAGonFS through
\texttt{/dev/fuse}. The file system keeps a \emph{namespace} of directories,
files and links, and a front-end buffer for each open file: the buffer is
loaded on \texttt{open} and stored on \texttt{flush} or \texttt{close}. A
\emph{data plane} moves block slices between ranks with non-blocking
point-to-point MPI messages and no collectives, so several ranks can transfer
at once. A \emph{control plane} sends one request message per rank. In
peer-to-peer mode a \emph{replayer} thread applies the namespace changes
made on other peers. Advisory locks (\texttt{fcntl}, \texttt{lockf},
\texttt{flock}) are served cluster-wide by one lock table on rank~0.

\begin{definitionbox}{Close-to-open consistency}
A write is acknowledged by every rank before \texttt{close} returns. A file
written and closed on one rank can therefore be read on any other rank as
soon as \texttt{close} returns. This matches DAGonStar's model exactly: a
consumer starts only after its producer has finished, and so after the
producer has closed its outputs.
\end{definitionbox}

\section{The Block Layout}

File contents are divided into blocks of $B = 128$\,KiB. A file of $S$ bytes
occupies $N = \lceil S/B \rceil$ blocks. Rank $r$ of $P$ owns
\[
c_r = \left\lfloor N/P \right\rfloor + [\,r < N \bmod P\,]
\]
consecutive blocks, starting at block $d_r = \sum_{q<r} c_q$. The layout
depends on $S$ and $P$ alone, so every rank can compute where any byte lives
without asking anyone. Lecture~A1 computes it for five sizes and checks four
invariants: every block has an owner, ranks differ by at most one block, rank~0
holds the most, and if rank~0 holds nothing, no rank does.

\lstinputlisting[caption={The DAGonFS block layout (\texttt{examples/A1-dagonfs/block\_layout.py}).},label={lst:block-layout}]{../../examples/A1-dagonfs/block_layout.py}

\lstinputlisting[style=bash,caption={Recorded output of \texttt{block\_layout.py}.}]{../../examples/A1-dagonfs/output/block_layout.txt}

A 564-byte file, the size of the course dataset, fits in one block on rank~0.
A 128\,MiB file is split into 1024 blocks, 256 per rank.

\section{A Workflow on a DAGonFS Volume}

DAGonStar and DAGonFS divide the work cleanly. DAGonStar schedules tasks,
resolves \texttt{workflow://} references, checkpoints and moves the final
data. DAGonFS answers POSIX calls inside its mount and never parses a
\texttt{workflow://} URI. The integration has four steps:
\begin{enumerate}[leftmargin=1.6em]
  \item start the volume before DAGonStar creates any task folder;
  \item set the workflow's \texttt{scratch\_dir\_base} to the mount point;
  \item run the workflow;
  \item copy the results to persistent storage (stage-out) before the volume
    is unmounted.
\end{enumerate}

The Python class \texttt{dagonfs.Volume(root, model, ranks, executable,
mpirun, mpirun\_args, local\_ranks, \ldots)} launches \texttt{mpirun} and
waits for the mount on \texttt{start}. It unmounts and waits for the ranks on
\texttt{stop}. Used as a context manager, it ties the life of the scratch
space to one block of the workflow script.

\lstinputlisting[caption={A DAGonStar workflow whose scratch lives on DAGonFS (\texttt{examples/A1-dagonfs/dagonfs\_workflow.py}).},label={lst:dagonfs-workflow}]{../../examples/A1-dagonfs/dagonfs_workflow.py}

\lstinputlisting[style=bash,caption={Recorded output of \texttt{dagonfs\_workflow.py} on a host without DAGonFS.}]{../../examples/A1-dagonfs/output/dagonfs_workflow.txt}

The recorded output comes from a laptop with no FUSE~3 and no DAGonFS build.
There the script prints the command it would launch, then runs the same
workflow on a temporary directory, so the workflow itself is always tested.
On a Linux node with MPI and FUSE~3, set \texttt{DAGONFS\_PYTHON} to the
\texttt{python} folder of the DAGonFS checkout and \texttt{DAGONFS} to the
executable, and the same script mounts a four-rank client-server volume.

\begin{cautionbox}[Stage out before unmounting]
Nothing survives the volume. A result that is needed after the workflow must
be copied to persistent storage inside the \texttt{with} block, as the
example does with \texttt{report.txt}. The last line of the output is read
\emph{after} the volume is gone, from the persistent copy.
\end{cautionbox}

\begin{sloppypar}
Mounting the volume first and pointing \texttt{scratch\_dir\_base} into it
puts the whole workflow on DAGonFS. Since version~3.0, DAGonStar also starts and
stops DAGonFS itself, through its built-in \texttt{dagonfs} ad-hoc provider:
\texttt{Workflow.create\_adhoc\_filesystem("dagonfs", ranks=4)} binds a volume
to selected tasks, and a \texttt{TaskType.VOLUME} task with
\texttt{provider="dagonfs"} makes the volume a node of the graph. The provider
supports the \texttt{fuse3}, \texttt{native} and \texttt{preload} frontends.
Appendix~\ref{app:volumes} describes both mechanisms.
\end{sloppypar}

\section{Building and Running}

DAGonFS requires CMake~3.20 or later, a C++23 compiler, MPI (Open MPI~4.1 or
later; \texttt{MPI\_THREAD\_MULTIPLE} for peer-to-peer) and libfuse~3.3 or
later.

\begin{lstlisting}[style=bash]
git clone https://github.com/DagOnStar/DAGonFS.git
cd DAGonFS
cmake -S . -B build -DUSE_MPI=ON -DCMAKE_BUILD_TYPE=Release
cmake --build build --parallel
ctest --test-dir build --output-on-failure

mpirun -np 4 build/dagonfs --model client-server -f /tmp/dagonfs
mpirun -np 4 build/dagonfs --model peer-to-peer  -f /tmp/dagonfs  # /tmp/dagonfs/0..3
fusermount3 -u /tmp/dagonfs                                      # stop the volume
\end{lstlisting}

\section{Performance}

The project compared DAGonFS with the NFS home directory of a cluster used as
the scratch space of real DAGonStar workflows (four ranks on one node,
November~2026). \Cref{tab:dagonfs-perf} reports the medians.

\begin{table}[h]
\centering
\begin{tabular}{@{}lrrr@{}}
\toprule
Workload (median, one node) & NFS & client-server & peer-to-peer \\
\midrule
File creation (ops/s) & 67 & 25\,532 & 26\,943 \\
4\,KiB files written (files/s) & 96 & 5\,923 & 7\,946 \\
DAGonStar, $8\times128$\,MiB, shell tasks (s) & 64.4 & 33.9 & 34.8 \\
DAGonStar, $8\times128$\,MiB, Python tasks (s) & 62.1 & 31.3 & 39.4 \\
DAGonStar, $2000\times8$\,KiB (s) & 257.6 & 24.7 & 28.2 \\
\bottomrule
\end{tabular}
\caption{DAGonFS against NFS as DAGonStar scratch space (from the DAGonFS
documentation).}
\label{tab:dagonfs-perf}
\end{table}

The large-file dataflow runs about 1.9 times faster than on NFS. The same
dataflow with Python tasks runs about 2.0 times faster, and the small-file
dataflow about 10.4 times faster. The gain is largest where it was expected:
many small files, whose cost on NFS is dominated by metadata operations.
Across two nodes DAGonFS stays 1.9 to 10.3 times faster. Unlike a node-local
\texttt{tmpfs}, it presents one namespace to tasks on different nodes, and
its capacity is the memory of the whole allocation.

\section{When to Use It}

\begin{itemize}[leftmargin=1.4em]
  \item Use DAGonFS for intermediate data that many tasks produce and
    consume within one allocation, especially many small files.
  \item Size the volume: the scratch data of the running tasks must fit in
    the memory of the ranks, alongside the tasks themselves.
  \item Do not use it for results that must outlive the job. Stage them out.
  \item Compare it with BeeOND (\Cref{ch:hpc-storage}). BeeOND uses the nodes'
    local disks and persists contents into task folders. DAGonFS uses memory
    and leaves stage-out to the workflow.
\end{itemize}

\begin{exercise}
Using the block-layout formula, compute $N$, $c_r$ and $d_r$ for a 1\,GiB file
on six ranks. Then estimate the memory each rank needs to hold the eight
128\,MiB intermediate files of the benchmark workflow at once.
\end{exercise}
